\section{Markov Chains}
\label{sec:markov_chain}
A \glsfirst{mc} \cite{levin2017markov} is a discrete-time stochastic process that describes the evolution of a state inside a space $\augmentedstatespace$ called the state space.
At each step, the probability of transition from a state $x\in\augmentedstatespace$ to a state $y\in\augmentedstatespace$ is given by the transition function $\markovchaintransition$ as $\markovchaintransition(y\vert x)$.

The properties of \glspl{mc} have been extensively studied in the literature, but are beyond the scope of this dissertation. 
We simply provide the following three definitions that will be useful in this work.
In the next definition, we denote by $X_t$ the random variable of the current state.


\begin{defi}[Recurrent state]
    A state $x\in\augmentedstatespace$ is recurrent if
    \begin{align*}
        \probability(\min\{t\ge1: X_t=x\}<\infty\vert X_0=x)=1
    \end{align*}
\end{defi}

\begin{defi}[Transient state]
    A state $x\in\augmentedstatespace$  that is not recurrent is transient.
\end{defi}

\begin{defi}[Unichain \gls{mc}]
\label{def:unichain_mc}
    A finite-state \gls{mc} is said to be unichain if it contains a single
    recurrent class, but can contain any number of transient states.
\end{defi}

The framework of \glspl{mc} encompasses many realistic processes. 
Yet, in some cases, one could wish to express the dependency on a set of older states for the current transitions. 
This is exactly what higher-order Markov chains--or Markov chains with memory--allow. 
In this framework, the transition can depend upon several past states.

\subsection{Higher-Order Markov Chains}
A \gls{mc} of order $l$ defines the transition to a state $x_t$ from the knowledge of the last $l$ states $(x_{t-l},\dots,x_{t-1})$ as,
\begin{align*}
    \probability(X_t = x_t\vert X_{t-l}=x_{t-l},\dots,X_{t-1}=x_{t-1})=P(x_t\vert x_{t-l},\dots,x_{t-1}).
\end{align*}

However, this model introduces many parameters to construct the transition probability.
Indeed, for a state space of size $\cardinal{\augmentedstatespace}=m$, the set $(x_{t-l},\dots,x_{t-1})$ can assume $m^l$ different values, and therefore probability $P$ needs $m^l(m-1)$ independent parameters to be defined \cite{berchtold2002mixture}.
Instead, \cite{raftery1985model} proposes a model with $m(m-1) + (l-1)$ independent parameters that retains a good modelling capacity in practice. 
This model, the \gls{mtd}, is presented below.
 
\begin{defi}[\Glsfirst{mtd} {\cite{raftery1985model}}]
\label{def:mixture_transition_distribution}
    An \gls{mtd} of order $l\in\naturalnumbers$ with state space $\augmentedstatespace=\{1,\dots,m\}$ is a \gls{mc} of order $l$ with the following properties.
    \begin{align*}
        \probability(X_t = x_t\vert X_{t-l}=x_{t-l},\dots,X_{t-1}=x_{t-1})
        &= \sum_{g=1}^l \lambda_g \probability(X_t=x_t\vert X_{t-1}=x_{t-g})
        \\
        &= \sum_{g=1}^l \lambda_g \markovchaintransition(x_t\vert x_{t-g}),
    \end{align*}
    where the probabilities $\markovchaintransition(x_t\vert x_{t-g})$ are defined as the transition of some \gls{mc} of order 1 and where
    \begin{align*}
        \sum_{g=1}^l \lambda_g &=1,
        \\
        \lambda_g &\ge 0.
    \end{align*}
\end{defi}

The model can be extended by introducing extra independent parameters to allow more modelling power.
This new model, with $lm(m-1) + (l-1)$ independent parameters, is given below.

\begin{defi}[\Glsfirst{mtdg} {\cite{raftery1985new,berchtold1996modelisation}}]
\label{def:multimatrix_mixture_transistion_distribution}
    The \gls{mtdg} model is an \gls{mtd} model with a modified transition matrix, 
    \begin{align*}
        \probability(X_t = x_t\vert X_{t-l}=x_{t-l},\dots,X_{t-1}=x_{t-1})
        &= \sum_{g=1}^l \lambda_g \probability(X_t=x_t\vert X_{t-g}=x_{t-g})
        \\
        &= \sum_{g=1}^l \lambda_g \markovchaintransition^{(g)}(x_t\vert x_{t-g}),
    \end{align*}
    where the transition probabilities $(\markovchaintransition^{(g)})_{g\in[\![1,l]\!]}$ come from the transition probabilities of $l$ different \glspl{mc} of order 1.
\end{defi}

However, it has been shown that the \gls{mtdg} model is over-parameterised \cite{lebre2008algorithm}.
Indeed, two different sets of parameters could define the same \gls{mtdg} model.
%\cite{lebre2008algorithm} propose a new %parameter set whose dimension is lower. It stems %from the straightforward remark that the %$l^{\text{nd}}$-order \gls{mtdg} model satisfies %the following proposition.
%
%\begin{prop}[{\cite[Proposition~1]{lebre2008algo%rithm}}]
%    Transition probabilities of a %$l^{\text{nd}}$-order \gls{mtdg} model %satisfy
%    \begin{align*}
%        \forall %&i_m,\dots,i_g,\dots,i_1,i_0,i_g'\in\sta%tespace
%        \\
%        &\probability(i_0\vert %i_m,\dots,i_g,\dots,i_1)-\probability(i_%0\vert i_m,\dots,i_g',\dots,i_1) = %\lambda_g \left[ q_{i_gi_0}^{(g)} - %q_{i_g'i_0}^{(g)} \right].
%    \end{align*}
%\end{prop}
%
%This proposition gives rives to a following %proposition.
%
%\begin{prop}[{\cite[Proposition~2]{lebre2008algo%rithm}}]
%    For any arbitrary $u$ element of %$\statespace$, the transition probabilities %of a mth-order \gls{mtdg} model satisfy:
%    \begin{align*}
%        &\forall i_m,\dots,i_0\in\statespace
%        \\
%        & \probability(i_0\vert %i_m,\dots,i_g,\dots,i_1) = \sum_{g=0}^l %\left[ p_u(i_0\vert i_g,g) - %\frac{l-1}{l} p_u(i_0)\right],
%    \end{align*}
%    where $p_u(j\vert i,g)=\probability(j\vert %u,\dots,i,\dots,u)$, with $u\dots uiu\dots %u$ is the $l$-letter word whose letter in %position $g$ (from right to left) is $i$ and %$p_u(j)$ is the value of $p_u(i\vert u,g)$, %whatever $g$.
%\end{prop}
%
%An \gls{mtdg} distribution can be parameterized %by a vector $\theta_u$ from the %$(m-1)[1+l(m-1)]$-dimensional
%set $\bar{\Theta}_u$,
%\begin{align*}
%    \bar{\Theta}_u =& \left\{(p_u(j\vert %i,g))_{0\le g\le l, i,j\in\statespace} \quad %\text{s.t.} \quad \forall g\in %\{1,\dots,l\}\forall i\in \statespace %\sum_{j\in\statespace}p(j\vert i,g)=1 %\right.
%    \\
%    & \left.\vphantom{\sum_{j\in\statespace}}\te%xt{ and } \forall g,g'\in \{1,\dots,l\}, %p_u(j\vert u,g)=p_u(j\vert u,g') \right\}.
%\end{align*}
%
%
%
%\begin{prop}[Inspired from %{\cite[Proposition~2]{lebre2008algorithm}}]
%    For arbitrary $u=(u_1,\dots,u_m)$ element of %$\statespace^m$, the transition %probabilities of a mth-order \gls{mtdg} %model satisfy:
%    \begin{align*}
%        &\forall i_m,\dots,i_0\in\statespace
%        \\
%        & \probability(i_0\vert %i_m,\dots,i_g,\dots,i_1) = \sum_{g=0}^l %\left[ p_u(i_0\vert i_g,g) - %\frac{l-1}{l} p_u(i_0)\right],
%    \end{align*}
%    where $p_u(j\vert i,g)=\probability(j\vert %u_0,\dots,i,\dots,u_m)$, with $u_0\dots %u_{g-1}iu_{g+1}\dots u_m$ is the $l$-letter %word whose letter in position $g$ (from %right to left) is $i$ and $p_u(j)$ is the %probability to reach $j$ from sequence $u$.
%\end{prop}


